\subsection{Mean-field limit of the self-training dynamics}
\label{subsec:mean-field}

Fix a finite iteration horizon $T\in\mathbb N$.
We describe the high-dimensional training dynamics through the
empirical distributions of coordinate trajectories.

Our main result constructs two effective stochastic processes,
$(\tilde w^t)_{t\geq0}$ and $(\tilde r^t)_{t\geq0}$,
through a coupled recursion.
Under the assumptions below, for every pair of
pseudo-Lipschitz functions of order $2$
(Definition~\ref{definition:pseudo-lipschitz})
$F,H:\mathbb R^{T+2}\to\mathbb R$, it establishes,
as $d\to\infty$,
\begin{align*}
    \frac1d\sum_{k=1}^d
        F(\mu_k,w_k^0,\ldots,w_k^T)
    &\xrightarrow{\mathbb P}
    \tilde{\mathbb E}
    \bigl[F(\tilde\mu,\tilde w^0,\ldots,\tilde w^T)\bigr],
    \\
    \frac1n\sum_{i=1}^n
        H(Y_i,r_i^0,\ldots,r_i^T)
    &\xrightarrow{\mathbb P}
    \tilde{\mathbb E}
    \bigl[H(\tilde Y,\tilde r^0,\ldots,\tilde r^T)\bigr].
\end{align*}
Here, $w^t\in\mathbb R^d$ and $b^t\in\mathbb R$ are generated
by Algorithm~\ref{alg:self-training}, and we define
\[
    r^t:=b^t\mathbf1_n+\frac{Xw^t}{\sqrt d},
    \qquad t\geq0.
\]
These formulas express limiting empirical averages as expectations
of the effective processes. Together with convergence of the
intercept, they determine the asymptotic signal alignment,
weight norm, and population classification error.

Intuitively, the effective processes describe the trajectories
of a typical weight coordinate and a typical training score,
much as an effective single-particle process describes the behavior
of a typical particle in a large interacting physical system.
Data reuse introduces memory terms and temporally correlated
Gaussian randomness, capturing the dependence on earlier
interactions with the training data.

The effective description depends on both the update functions
and the asymptotic behavior of the initialization.
We first impose regularity conditions on the loss and penalty.

\begin{assumption}[Regularity of loss and penalty]
\label{assumption:regularity-of-loss-and-penalty}
The functions $r\mapsto\ell(y,r)$, for $y\in\{-1,+1\}$,
and $u\mapsto j(u)$ are differentiable.
There exist dimension-independent constants
$L_\ell,L_j\in(0,\infty)$ such that
\[
    |\ell'(y,r)-\ell'(y,s)|
    \leq L_\ell|r-s|,
    \qquad
    |j'(u)-j'(v)|
    \leq L_j|u-v|
\]
for every $y\in\{-1,+1\}$ and $r,s,u,v\in\mathbb R$.
Moreover, $\ell'(y,\cdot)$ is nonconstant for each
$y\in\{-1,+1\}$.
\end{assumption}
The initial parameters $b^{\mathrm{init}}\in\mathbb R$ and
$w^{\mathrm{init}}\in\mathbb R^d$ are deterministic and may
depend on $d$. Define the joint empirical law of the signal
and initial weight coordinates by
\[
    \nu_{(\mu,w^{\mathrm{init}})}
    :=
    \frac1d\sum_{k=1}^d
        \delta_{(\mu_k,w_k^{\mathrm{init}})}.
\]
We also impose the following asymptotic conditions.

\begin{assumption}[Quenched initialization]
\label{assumption:quenched-initialization}
The following conditions hold as $d\to\infty$:
\begin{enumerate}
    \item Convergence of the joint empirical law:
    \begin{equation}
        \nu_{(\mu,w^{\mathrm{init}})}
        \xrightarrow{W_2}
        \bar\nu_{(\mu,w^{\mathrm{init}})}
        \in\mathscr P_2(\mathbb R^2).
        \label{eq:backward-init}
    \end{equation}

    \item Convergence of the initial intercept:
    \begin{align}
        b^{\mathrm{init}}
        \longrightarrow
        \bar b^{\mathrm{init}}\in\mathbb R.
    \end{align}

    \item Positive limiting weight second moment:
    \begin{equation}
        (\bar\tau^{\mathrm{init}})^2
        :=
        \int_{\mathbb R^2}w^2\,
        \bar\nu_{(\mu,w^{\mathrm{init}})}(du,dw)
        >0.
        \label{eq:positive-init-mean-field-weight-norm}
    \end{equation}
\end{enumerate}
\end{assumption}
By Lemma~\ref{lemma:wasserstein-lipschitz-pushforward}
and uniqueness of Wasserstein limits, $\bar\nu_\mu$ is the
first marginal of $\bar\nu_{(\mu,w^{\mathrm{init}})}$.
The limiting law $\bar\nu_{(\mu,w^{\mathrm{init}})}$ and intercept
$\bar b^{\mathrm{init}}$ provide the initial data for the effective
description. By requiring the initialization to be deterministic
relative to the Gaussian design noise and to have a positive
limiting weight second moment, we exclude zero initialization
and classifiers fitted to the same training data, while allowing
classifiers trained exogenously. Throughout this subsection,
we also fix the hyperparameters $\eta>0$, $\lambda,\pi\geq0$,
and $\kappa_-<0<\kappa_+$ independently of the dimension.

\subsubsection{Construction of the effective processes}

The effective processes are constructed on two probability spaces. Let
\[
    (\tilde\Omega^{[g]},\tilde{\mathcal F}^{[g]},
        \tilde{\mathbb P}^{[g]})
    \qquad\text{and}\qquad
    (\tilde\Omega^{[w]},\tilde{\mathcal F}^{[w]},
        \tilde{\mathbb P}^{[w]})
\]
be, respectively, the sample-side and parameter-side mean-field
probability spaces. The sample-side space supports
\[
    (\tilde Y,\tilde\Delta)\sim\bar\nu_{(Y,\Delta)}
    \qquad\text{and}\qquad
    (z_t^{[g]})_{t\geq0}
    \stackrel{\mathrm{i.i.d.}}{\sim}\mathcal N(0,1),
\]
while the parameter-side space supports
\[
    (\tilde\mu,\tilde w^0)
    \sim\bar\nu_{(\mu,w^{\mathrm{init}})}
    \qquad\text{and}\qquad
    (z_t^{[w]})_{t\geq0}
    \stackrel{\mathrm{i.i.d.}}{\sim}\mathcal N(0,1).
\]
On each space, the Gaussian sequence is independent of the
corresponding pair of base variables.
The two recursions interact through deterministic expectations
and coefficients; no joint law between their random variables
is prescribed. The convergence theorem compares their laws
with the finite-dimensional empirical measures, without requiring
a coupling with the original training dynamics.

In the following, we write simply $\tilde{\mathbb P}$ and
$\tilde{\mathbb E}$ for probability and expectation on either
mean-field probability space. The relevant space is determined by
the random variables appearing in the expression. For notational convenience, we also set $\ell'(0,r):=0$ for every
$r\in\mathbb R$, so that unselected observations contribute zero.

\begin{definition}[Mean-field recursion]
\label{def:mean-field-recursion}
Fix a deterministic initial intercept $\bar b^0:=\bar b^{\mathrm{init}}$. We recursively define 
\begin{align*}
    (\tilde q^t,\tilde r^t,\tilde{\mathfrak c}^t,\tilde g^t)_{t\ge0}, 
    \quad 
    (\tilde p^t,\tilde w^t)_{t\ge0},
    \quad
    (\bar b^t,\bar m^t,\bar\zeta^t,\bar\chi^t)_{t\ge0},
\end{align*}
the sample-side processes, the parameter-side processes, and the deterministic macroscopic quantities respectively. We also define the projection, memory, and innovation coefficients below.
Write the weight and residual correlation kernels as
\begin{align}
    \bar C^{[w]}(s,t)
    &:=
    \tilde{\mathbb E}
    [
        \tilde w^s\tilde w^t
    ],
    &
    \bar C^{[g]}(s,t)
    &:=
    \tilde{\mathbb E}
    [
        \tilde g^s\tilde g^t
    ].
    \label{eq:correlation-kernel-definition}
\end{align}
Under the hypotheses of Theorem~\ref{theorem:trajectory},
Lemma~\ref{lemma:mf-proof-nondegeneracy} establishes
well-definedness of the recursion, uniqueness of the
coefficient systems' solutions, and finite second moments
of all coordinates. It also establishes absolute continuity
of the score law at every finite time.

The construction proceeds in the order
\[
    \tilde w^t
    \longrightarrow(\tilde q^t,\tilde r^t,\tilde g^t)
    \longrightarrow(\tilde p^t,\tilde w^{t+1}).
\]

For every $t\ge0$, define the following quantities:

\begin{itemize}[wide]
    \item Weight-side projection and memory coefficients:
    \begin{align}
    \sum_{s=0}^{t-1}
        \bar C^{[w]}(u,s)\bar\alpha^{[w]}(t,s)
    &=
    \bar C^{[w]}(u,t),
    \qquad 0\le u<t,
    \label{eq:alpha-w-definition}
    \\[-7pt]
    \tilde w_\perp^t
    &:=
    \tilde w^t
    -
    \sum_{s=0}^{t-1}
        \bar\alpha^{[w]}(t,s)\tilde w^s,
    \\[-7pt]
    \sum_{s=0}^{t-1}
        \bar C^{[g]}(u,s)\bar\phi^{[w]}(t,s)
    &=
    \tilde{\mathbb E}
    [
        \tilde p^u\tilde w_\perp^t
    ],
    \qquad 0\le u<t,
    \label{eq:phi-w-definition}
    \\[-7pt]
    \bar v_t^{[w]}
    &:=
    \left(
        \tilde{\mathbb E}
        [
            (\tilde w_\perp^t)^2
        ]
    \right)^{1/2}.
\end{align}

    \item Residual-side projection and memory coefficients:
    \begin{align}
    \sum_{s=0}^{t-1}
        \bar C^{[g]}(u,s)\bar\alpha^{[g]}(t,s)
    &=
    \bar C^{[g]}(u,t),
    \qquad 0\le u<t,
    \label{eq:alpha-g-definition}
    \\[-7pt]
    \tilde g_\perp^t
    &:=
    \tilde g^t
    -
    \sum_{s=0}^{t-1}
        \bar\alpha^{[g]}(t,s)\tilde g^s,
    \\[-7pt]
    \sum_{s=0}^{t}
        \bar C^{[w]}(u,s)\bar\phi^{[g]}(t,s)
    &=
    \tilde{\mathbb E}
    [
        \tilde q^u\tilde g_\perp^t
    ],
    \qquad 0\le u\le t,
    \label{eq:phi-g-definition}
    \\[-7pt]
    \bar v_t^{[g]}
    &:=
    \left(
        \tilde{\mathbb E}
        [
            (\tilde g_\perp^t)^2
        ]
    \right)^{1/2}.
\end{align}

    \item Cumulative forward memory coefficients and Gaussian process:
    \begin{align}
    \bar\gamma^{[w]}(t,s)
    &:=
    \bar\phi^{[w]}(t,s)
    +
    \sum_{u=s+1}^{t-1}
        \bar\alpha^{[w]}(t,u)
        \bar\gamma^{[w]}(u,s),
    \qquad 0\le s<t,
    \\[-7pt]
    \xi_t^{[g]}
    &:=
    \sum_{u=0}^{t-1}
        \bar\alpha^{[w]}(t,u)\xi_u^{[g]}
    +
    \bar v_t^{[w]}z_t^{[g]}.
    \label{eq:xi-g-definition}
\end{align}

    \item Cumulative backward memory coefficients and Gaussian process:
\begin{align}
    \bar\gamma^{[g]}(t,s)
    &:=
    \bar\phi^{[g]}(t,s)
    +
    \sum_{u=s}^{t-1}
        \bar\alpha^{[g]}(t,u)
        \bar\gamma^{[g]}(u,s),
    \qquad 0\le s\le t,
    \label{eq:Gamma-g-definition}
    \\[-7pt]
    \xi_t^{[w]}
    &:=
    \sum_{u=0}^{t-1}
        \bar\alpha^{[g]}(t,u)\xi_u^{[w]}
    +
    \bar v_t^{[g]}z_t^{[w]}.
    \label{eq:xi-w-definition}
\end{align}

\item Forward fluctuation:
\vspace{-21pt}
\begin{align}
    \tilde q^t
    &:=
    \frac{1}{\sqrt{\bar\delta}}
    \sum_{s=0}^{t-1}
        \bar\gamma^{[w]}(t,s)\tilde g^s
    +
    \xi_t^{[g]}.
    \label{eq:q-memory-definition}
\end{align}

\item Backward fluctuation:
\vspace{-21pt}
\begin{align}
    \tilde p^t
    &:=
    \sqrt{\bar\delta}
    \sum_{s=0}^{t}
        \bar\gamma^{[g]}(t,s)\tilde w^s
    +
    \xi_t^{[w]}.
    \label{eq:p-memory-definition}
\end{align}

\item The macroscopic quantities:
\begin{align}
    \bar b^{t+1}
    :=
    \bar b^t
    -
    \eta\bar\zeta^t,
    \qquad
    \bar m^t
    :=
    \tilde{\mathbb E}
    [
        \tilde\mu\tilde w^t
    ],
    \qquad
    \bar\zeta^t
    :=
    \tilde{\mathbb E}
    [
        \tilde g^t
    ],
    \qquad
    \bar\chi^t
    :=
    \tilde{\mathbb E}
    [
        \tilde Y\tilde g^t
    ].
\end{align}

\item The sample side process:
\begin{align}
    \tilde r^t
    :=
    \bar b^t
    +
    \bar m^t\tilde Y
    +
    \sigma\tilde q^t,
    \qquad
    \tilde{\mathfrak c}^t
    :=
    \mathbf 1_{\{\tilde r^t\ge\kappa_+\}}
    -
    \mathbf 1_{\{\tilde r^t\le\kappa_-\}}.
    \label{eq:mean-field-score}
\end{align}
\vspace{-21pt}
\begin{align}
    \tilde g^t
    :=
    \tilde g_{\mathcal L}^t
    +
    \pi\tilde g_{\mathcal U}^t,
    \qquad
    \tilde g_{\mathcal L}^t
    :=
    \frac{\tilde\Delta}{\bar\rho}
    \ell'
    \left(
        \tilde Y,\tilde r^t
    \right),
    \qquad
    \tilde g_{\mathcal U}^t
    :=
    \frac{1-\tilde\Delta}{1-\bar\rho}
    \left| \tilde{\mathfrak c}^t \right|
    \ell'
    \left(
        \tilde{\mathfrak c}^t,\tilde r^t
    \right).
    \label{eq:g-U-definition}
\end{align}

\item The parameter side process:
\begin{align}
    \tilde w^{t+1}
    &:=
    \tilde w^t
    -
    \eta
    \left[
        \lambda j'(\tilde w^t)
        +
        \bar\chi^t\tilde\mu
        +
        \frac{\sigma}{\sqrt{\bar\delta}}
        \tilde p^t
    \right].
    \label{eq:w-update-definition}
\end{align}
\end{itemize}
\end{definition}

\subsubsection{Mean-field convergence}

The following theorem identifies the effective trajectory laws
as the limits of the finite-dimensional empirical measures.

\begin{theorem}[Mean-field trajectory convergence]
\label{theorem:trajectory}
Suppose that Assumption~\ref{assumption:asymptotic-regime},
Assumption~\ref{assumption:regularity-of-loss-and-penalty},
and Assumption~\ref{assumption:quenched-initialization} hold.

For every fixed $T\in\mathbb N$, the following convergences hold
as $d\to\infty$:
\begin{enumerate}[label=(\roman*),wide]
    \item \emph{Parameter-side trajectory convergence.}
    \begin{align}
        W_2\left(
            \frac1d\sum_{k=1}^d
                \delta_{(\mu_k,w_k^0,\ldots,w_k^T)},
            \mathscr L(\tilde\mu,\tilde w^0,\ldots,\tilde w^T)
        \right)
        &\xrightarrow{\mathbb P}0.
        \label{eq:parameter-trajectory-limit}
    \end{align}

    \item \emph{Sample-side trajectory convergence.}
    \begin{align}
        W_2\left(
            \frac1n\sum_{i=1}^n
                \delta_{(Y_i,\Delta_i,r_i^0,\ldots,r_i^T)},
            \mathscr L(
                \tilde Y,\tilde\Delta,\tilde r^0,\ldots,\tilde r^T
            )
        \right)
        &\xrightarrow{\mathbb P}0.
        \label{eq:sample-trajectory-limit}
    \end{align}

    \item \emph{Intercept convergence.}
    \begin{align}
        \max_{0\leq t\leq T}|b^t-\bar b^t|
        &\xrightarrow{\mathbb P}0.
        \label{eq:mean-field-bias-convergence}
    \end{align}
\end{enumerate}
Here, the finite dimensional sequences $(w^t)_{t \ge 0}$, $(b^t)_{t \ge 0}$ and $(r^t)_{t \ge 0}$ are generated by Algorithm~\ref{alg:self-training}, and
probability is taken with respect to the finite-dimensional
Gaussian design noise. The limiting trajectory laws are
deterministic probability measures.
\end{theorem}
By Corollary~\ref{corollary:wasserstein-empirical-observables},
the trajectory convergence implies convergence in probability
of empirical averages of pseudo-Lipschitz functions of order $2$.
The following corollary records the resulting limits for
macroscopic observables and classification error.
Define
\begin{align}
    m^t
    &:=
    \frac{\mu^\top w^t}{d},
    &
    \tau^t
    &:=
    \frac{\|w^t\|}{\sqrt d},
    \\
    C^{[w]}(s,t)
    &:=
    \frac{(w^s)^\top w^t}{d},
    &
    \omega^t
    &:=
    \frac1M\sum_{j\in I_{\mathcal U}}S_\kappa(r_j^t).
    \label{eq:finite-dimensional-trajectory-observables}
\end{align}
The quantity $m^t$ measures signal alignment, $\tau^t$ is the normalized
weight norm, and $\omega^t$ is the selected fraction of unlabeled
observations.
Define also the selected-error fraction
\[
    e^t
    :=
    \frac1M\sum_{j\in I_{\mathcal U}}
        \mathbf1_{\{\mathfrak C(\theta^t)_j=-Y_j\}}.
\]
This is the fraction of all unlabeled observations that
are selected with an incorrect pseudo-label.
The corresponding limiting weight norm is
\[
    \bar\tau^t
    :=
    \left(\tilde{\mathbb E}[(\tilde w^t)^2]\right)^{1/2}
    =
    \sqrt{\bar C^{[w]}(t,t)}
    >0.
\]
Positivity follows from
Lemma~\ref{lemma:mf-proof-nondegeneracy}.

\begin{corollary}[Macroscopic observables and classification error]
\label{cor:mean-field-observables}
Under the assumptions of Theorem~\ref{theorem:trajectory}, for every
fixed $s,t\geq0$,
\begin{align}
    m^t
    &\xrightarrow{\mathbb P}\bar m^t,
    &
    \tau^t
    &\xrightarrow{\mathbb P}\bar\tau^t,
    &
    C^{[w]}(s,t)
    &\xrightarrow{\mathbb P}\bar C^{[w]}(s,t),
    &
    \omega^t
    &\xrightarrow{\mathbb P}\bar\omega^t,
    &
    e^t
    &\xrightarrow{\mathbb P}\bar e^t,
    \label{eq:mean-field-observable-convergence}
\end{align}
where
\[
    \bar\omega^t
    :=
    \frac{
        \tilde{\mathbb E}
        [(1-\tilde\Delta)S_\kappa(\tilde r^t)]
    }{1-\bar\rho},
    \qquad
    \bar e^t
    :=\tilde{\mathbb P}
        (\tilde{\mathfrak c}^{\,t}=-\tilde Y\mid\tilde\Delta=0).
\]

Let $\mathcal D$ denote the observed training data.
Independently of $\mathcal D$, draw
$Y_{\mathrm{test}}\in\{-1,+1\}$ with
$\mathbb P(Y_{\mathrm{test}}=+1)=p$ and
$Z_{\mathrm{test}}\sim\mathcal N(0_d,I_d)$, independently, and set
\[
    X_{\mathrm{test}}
    :=\frac{Y_{\mathrm{test}}\mu}{\sqrt d}+\sigma Z_{\mathrm{test}}.
\]
Define the conditional test error by
\[
    \mathcal E^t
    :=\mathbb P\!\left(
        f_{\theta^t}(X_{\mathrm{test}})\neq Y_{\mathrm{test}}
        \,\middle|\,\mathcal D
    \right).
\]
Then, for every fixed $t\geq0$,
\begin{align}
    \mathcal E^t
    &\xrightarrow{\mathbb P}
    \bar{\mathcal E}^t
    :=
    \bar p\,
    \Phi\left(
        \frac{-\bar b^t-\bar m^t}{\sigma\bar\tau^t}
    \right)
    +
    (1-\bar p)\,
    \Phi\left(
        \frac{\bar b^t-\bar m^t}{\sigma\bar\tau^t}
    \right),
    \label{eq:mean-field-classification-error}
\end{align}
where $\Phi$ is the standard Gaussian distribution function.
\end{corollary}

The proof of Theorem~\ref{theorem:trajectory} uses iterative
Gaussian conditioning following
Bolthausen and Bayati--Montanari
\cite{bolthausen2014,bayati2011dynamics};
it is presented in Appendix~\ref{sec:proof-mean-field}.
We apply conditioning directly to the alternating forward
and backward matrix products, deriving memory coefficients
through orthogonal projections rather than a
reduction to AMP. 
Defining these coefficients through Gaussian inner products
avoids differentiating the discontinuous selection rule and
retains the contributions associated with its threshold jumps.

The main probabilistic step is to reveal the history through
successive linear observations of the Gaussian design matrix.
Although the observation directions depend on previous
iterates, they are fixed conditional on the preceding history.
The conditional matrix law therefore remains Gaussian and
admits a decomposition into its conditional mean and a
projected Gaussian term generated from fresh randomness.
The conditional means produce the memory terms, while the
remaining randomness produces the innovations.

The rest of the argument establishes convergence of the
projection coefficients and empirical trajectory laws.
In particular, it verifies nondegeneracy and controls the
discontinuous feedback using the absence of atoms at its
discontinuity points.
